% Hafta 4 — Yığın (stack) ve öbek (heap)
% Derleme: py -3.12 tools/latex/derle.py docs/week-4/assets/h04-13-yigin-obek.tex
\documentclass[tikz,border=8pt]{standalone}
\usepackage{cen429-cizim}
\begin{document}
\begin{tikzpicture}[node distance=10mm and 24mm]
  \node[baslik] (b) at (0,0) {Yığın (stack) ve öbek (heap)};
  \node[altbaslik] at (b.south west) {taşmanın sonucu, belleğin hangi bölgesinde olduğuna göre değişir};

  \node[kutu=58mm, below=16mm of b.south west, anchor=north west] (yigin) {%
    \kb{YIĞIN}\\[2mm]
    fonksiyonla açılır/kapanır\\
    otomatik, hızlı\\
    boyut derleme anında belli\\[3mm]
    taşma $\rightarrow$ dönüş adresi tehlikede};
  \node[uyarikutu=58mm, right=of yigin] (obek) {%
    \kb{ÖBEK}\\[2mm]
    malloc/free ile yönetilir\\
    yaşam süresi size bağlı\\
    boyut çalışma anında\\[3mm]
    taşma $\rightarrow$ komşu blok / meta veri};

  \node[fit=(yigin)(obek), inner sep=0pt] (grup) {};
  \node[dip, text width=150mm, below=10mm of grup.south, anchor=north] {Kanarya yığını korur; öbek taşmasını yakalamaz --- bu yüzden ASan gerekir.};
\end{tikzpicture}
\end{document}
